\scanpage{241}
\mspage{123}

\deleted{la subtilité que \inserted{voudrai} \unsure{j'introduirai|j'introduirais} dans mon groupe $G$}

Je considère cet $H$ comme une structure supplémentaire, en plus du \deleted{$Z$} groupe d'homothéties $Z$, et définirai $G$ en conséquence par

\[
\begin{aligned}
G={}&\deleted{\text{groupe des}}\ \bigl\{f\in\unsure{\mathfrak S|\mathcal G}_S\mid f\text{ normalise }G,\text{ centralise}\\
& Z,\ \text{et }\exists\lambda\in Z\text{ tel}\\
&\text{que } f(A_0)=A_\lambda\quad\text{i.e.}\\
&\forall s,t\in S,\quad\{s,t\}\in A_0\Longleftrightarrow\{f^{-1}s,\unsure{f\lambda^{-1}t|f^{-1}\lambda t}\}\in A_0
\end{aligned}
\tag*{(63)}
\]

de sorte qu'on aura \deleted{\unsure{encore}} un diagramme

\diagram{\begin{tikzcd}[ampersand replacement=\&]
\&\&H\arrow[d,hook]\arrow[dr,dashed,"\text{\uline{distingué}}"]\&\&\\
{(64)}\&1\arrow[r]\&\underset{\substack{\shortparallel\\\mathrm{Ker}\,\delta}}{G_0}\arrow[r]\&G\arrow[r,"\delta",no head]\&Z\\
\&\&\&Z\arrow[u,hook]\arrow[ur,"{\lambda\mapsto\lambda^2}"']\&
\end{tikzcd}}

\margin{$\nearrow$ Pour écrire ici $Z$, et non un quotient de $Z$, on admet que $Z$ opère \uline{fidèlement} sur $\pi$, ce qui est automatique si tt sommet de $\gamma_0$ est déterminé par son étoile}

où l'inclusion composée $H\longrightarrow G_0\unsure{-|\to}G$ fait de $H$ un sous-groupe \uline{distingué} de $G$. Pour l'instant on ne suppose pas $\delta$ surjectif. On fera attention que dans cette présentation, même si on suppose $\delta$ surjectif, les groupes $G_0$, \inserted{$G$} introduits ici \deleted{\unsure{ne sont pas}} \inserted{sont à priori} plus petits que ceux du n$^\circ$ précédent, qui eux ne font pas intervenir la donnée particulière \unsure{du hom|du choix|d'un} $H$. En fait, ceux \deleted{derniers \unsure{deux}} \inserted{envisagés ici sont} les \uline{normalisateurs} de $H$ dans les groupes $G_0^{\unsure{!}},G$ envisagés au n$^\circ$ précédent. L'on retrouvera \deleted{la situation} ces derniers, bien sûr, en prenant $H=G_0$.

Notre objet principal est de décrire de façon aussi précise que possible le diagramme (64),

% archive sheet number: 240
